% RESEARCH_DRAFT_SHA256: 2eaf2337a2fd8e86f8771e0e8d9776d93b1e061ce703070c44091678d64ee001
% Layout and source-reference wrappers expanded for the web renderer.
\documentclass[11pt]{article}
\usepackage[margin=25mm]{geometry}
\usepackage{amsthm}
\usepackage{fontspec,unicode-math}
\setmainfont{Latin Modern Roman}
\setmathfont{Latin Modern Math}
\usepackage{amsmath,amssymb,mathtools,xurl,hyperref}
\hypersetup{colorlinks=true,urlcolor=blue,linkcolor=blue}
\setcounter{secnumdepth}{0}
\setlength{\parindent}{0pt}
\setlength{\parskip}{5pt}
\setlength{\emergencystretch}{3em}
\newcommand{\sref}[2]{\hyperlink{source:#1:#2}{[S#2]}}
\newcommand{\cataloguescope}[1]{\paragraph{Source record.}#1}
\newcommand{\floor}[1]{\left\lfloor#1\right\rfloor}
\newtheorem{proposition}{Proposition}
\newtheorem{lemma}[proposition]{Lemma}
\begin{document}
\paragraph{Research manuscript.} OpenAI \texttt{gpt-6-astra}, reasoning \texttt{ultra}; 11 October 2026.
\paragraph{Preserved source body.} The complete catalogue mathematical text and sources follow with layout and source links expanded, including source qualifications and its historical preparation notes. New research begins at the section headed Research attempt.

\section{A versioned closed-loop audit for STACK}\label{statement}

\subsection{Definitions and mathematical statement}

\paragraph{Shared notation.}Standard number systems are $\mathbb N,\mathbb Z,\mathbb Q,\mathbb R,\mathbb C$. Specialized notation and index ranges follow the cited source and the definitions below.


Instruction comprises teaching activities and decisions. An assessment tool can change those activities through available tasks, feedback, and observed responses. Evaluating reinforcement requires an explicit account of the intended instruction and a measured outcome.
\paragraph{Record scope and qualifications.}
The catalogue title is retained as metadata. Named certificates, contracts, bounds or frameworks in that title are not asserted as results of the original source extract.
The prompt gives no course, intervention, measurement, or comparison. It remains an educational research agenda. The supplied catalogue marks this record partially solved; its present research status was not independently checked.
The source text below is rendered literally where necessary; damaged formulas and composite extracts are not silently repaired into a new mathematical claim.
\paragraph{Statement recorded in the supplied source.}
19. How does Stack change instruction? How does Stack reinforce instruction? \href{http://aimath.org/pastworkshops/oerassessmentproblems.pdf}{[S2]}
\paragraph{Source status.}
The supplied catalogue labels this source record partially solved; that metadata was checked against the supplied snapshot on 2026-10-11. The label is preserved for provenance and is not an independent verification of present open status or mathematical correctness.
\subsection{Short English statement}
The source asks how STACK changes or reinforces instruction.
\subsection{Sources}
\begin{itemize}
\item[S1] Ulam AI and the UnsolvedMath Contributors, UnsolvedMath: A Curated Collection of Open Mathematics Problems, catalogue entry. Catalogue curation and metadata: CC BY 4.0. The source extract is retained; definitions, English summaries and qualifications are editorial additions.
\url{https://huggingface.co/datasets/ulamai/UnsolvedMath}.
\item[S2] AIM Workshop Problem Lists, source linked by the catalogue.
\url{http://aimath.org/pastworkshops/oerassessmentproblems.pdf}.
\end{itemize}
\paragraph{Preparation and review.}
The original source and UnsolvedMath contributors retain attribution for the source extract. Codex (GPT-6; exact serving revision not exposed) assisted with definitions, English summaries, scope notes and formatting. No new human mathematical authorship or independently verified solution is asserted. Catalogue identity checked against the complete supplied registry and current repository. Distinct upstream records, including repeated questions and broad or damaged extracts, are retained for complete coverage of the source CS category. A repeated record is not a new independent conjecture. Complete mathematical-equivalence and current-literature review has not been performed. Preparation and final review follow contribution prompts C and D, with source gaps disclosed. Structural and rendering checks do not establish a new theorem or open status.
\label{end:statement}

\section{Research attempt}
\subsection{Literature, source context and scope}
The primary AIM notes describe a brainstorming session. The original excerpt and all catalogue qualifications remain above. They do not specify a mathematical theorem or a complete empirical study. The formal model below is an explicitly limited research route introduced for this attempt, not a reconstruction of a missing original conjecture. The supplied catalogue's earlier research summary was inspected as context; its named framework is not adopted as an established solution. The elementary result below is not claimed novel, and no educational effectiveness data were collected.

\subsection{Attempt: a delayed instructional-response loop}
To test the inference that a stronger corrective response always reinforces a process, consider a declared linear controller for an error variable $e_t$. Immediate correction follows $e_{t+1}=(1-k)e_t$. A one-step delay instead follows $e_{t+1}=e_t-ke_{t-1}$. This is an illustrative control model, not an empirical law for learners or a description of STACK.
\begin{proposition}
For $k>0$, every immediate-correction trajectory tends to zero exactly when $k<2$. Every delayed-correction trajectory tends to zero exactly when $k<1$. Hence a correction gain can stabilize the immediate loop and destabilize the delayed loop.
\end{proposition}
\begin{proof}
The immediate solution is $e_t=(1-k)^te_0$, so convergence for every initial value is equivalent to $|1-k|<1$. For the delayed recurrence the characteristic polynomial is $r^2-r+k$. When $0<k<1/4$, its roots $(1\pm\sqrt{1-4k})/2$ are both strictly between zero and one, so every linear combination of their powers tends to zero. At $k=1/4$, the repeated root is $1/2$ and solutions are $(A+Bt)2^{-t}$, also converging to zero. For $k>1/4$, the conjugate roots have modulus $\sqrt{k}$ because their product is $k$. Thus all solutions converge when $k<1$, whereas at $k=1$ nonzero periodic solutions exist and for $k>1$ some real solutions have unbounded subsequences. To see the last assertion, write a nonzero solution as $k^{t/2}\cos(t\theta+\phi)$; the pair of consecutive cosine factors cannot both tend to zero when $\sin\theta\ne0$. This proves the stated all-trajectory criterion.
\end{proof}
\subsection{Verification and remaining gaps}
At $k=3/2$, the immediate multiplier is $-1/2$, but delayed roots have modulus $\sqrt{3/2}>1$. Iteration from $e_0=e_1=1$ checks the divergent oscillatory behavior over a finite horizon; the root argument supplies the unrestricted proof. The failed inference is retaining an immediate-response stability rule after adding a delay. Whether an educational deployment has any such state, gain or delay relation requires evidence. No instructional effect or optimal feedback gain is inferred from this controller example.

\section{References and source attribution}
American Institute of Mathematics, workshop discussion notes on open educational resources and assessment, \url{https://aimath.org/pastworkshops/oerassessmentproblems.pdf}. The 16-page primary PDF was opened on 11 October 2026 and its numbered brainstorming item was checked against the supplied record. The source supports the provenance and agenda, not the additional modeling assumptions or an empirical educational conclusion. Catalogue and contributor attribution are retained in the reproduced Sources section. The worked mathematics is proved here without reliance on an unverified external theorem.


\paragraph{Integration review.} Codex (GPT-6; exact serving revision not exposed) reviewed the written arguments and expanded layout and source-reference wrappers. Repository metadata and the production web renderer were checked using Python 3.13 and JavaScript with MathJax 3.2.2. This model-assisted review is not independent human proof verification; no PDF compilation or formal proof checking is asserted.

\section{Completion Estimate}
\noindent\textbf{COMPLETION: 0\%}
No substantiated new progress toward the entire original target is claimed. The worked results and failed route are delimited above; a known result, illustrative model, or auxiliary obstruction is not counted as original-target completion.

\section{Final status and attribution}
\textbf{Attempt status: unresolved.} This describes this research attempt, not a new ruling on the source catalogue's status or on external literature claims.
The mathematics and this self-audit were generated by OpenAI \texttt{gpt-6-astra}, reasoning \texttt{ultra}, on 11 October 2026. Codex (GPT-6; exact serving revision not exposed) orchestrated the batch and contributed formatting/repository review. No human mathematical authorship, independent human proof checking, or formal verification is claimed. The historical preparation paragraphs in the reproduced catalogue body concern that body only.

Contribution prompt A was applied to the complete definition and supplied source record, followed by the adversarial checks of prompt B. Local reproduction files are \texttt{batch008/checks.py}, \texttt{computations.json}, and the per-ID audit JSON. They were not public links at preparation time. Numerical checks are finite and do not prove unrestricted results. No LaTeX engine was found locally; compilation was not performed by this research agent.

\end{document}
